\chapter{Lab 4}


\section{评分标准}
\textbf{60\%验收+40\%报告}

验收部分八个题,每个同学问两个总共占验收分数的30\%，70\% 实验部分由 35\% 仿真和 35\% 上板组成。

报告部分分数由思考题组成，共10个思考题，每个思考题若回答不正确扣5分，若回答不完全正确扣1-4分。


\section{预设问题}

\begin{enumerate}
    \item \textbf{Q：}
    为什么在链接脚本中，需要将 sp 指向 \_ekernel 后 4 KiB 的位置，将其作为栈顶？

    \textbf{A：}
    能根据内存布局，讲出防止在后续过程中修改 .bss 等数据段即可。

    \item \textbf{Q：}
    SBI是什么？我们为什么需要SBI？OpenSBI是什么？

    \textbf{A：}
    SBI（Supervisor Binary Interface）是 S-mode 的 kernel 和 M-mode 执行环境之间的接口规范，而 OpenSBI 是一个 RISC-V SBI 规范的开源实现。需要SBI的原因是操作系统内核需要操作硬件，但是本身权限不够，需要借助SBI接口获得部分M态权限。

    \item \textbf{Q:}
    C的内联汇编语法有四个部分组成，分别是哪四个部分？

    \textbf{A:}
    第一部分是汇编指令，指令末尾需要添加换行符或者 ; 作为指令之间的分隔符。
    第二部分是输出操作数部分。
    第三部分是输入操作数部分。
    第四部分是可能影响的寄存器或存储器，用于告知编译器当前内联汇编语句可能会对某些寄存器或内存进行修改，使得编译器在优化时将其因素考虑进去。

    \item \textbf{Q:}
    简述如何通过代码插桩或gdb调试判断当前特权态

    \textbf{A:}
    尝试读取S态/M态寄存器，有以下结论：
    都不能读取：Umode；只能读S态：Smode；S/M态寄存器都可读：Mmode

    

    

    

\end{enumerate}